% !TEX root = ../ConvexOptimization.tex

\chapter{Convex sets}

% -- BEGIN LECTURE 2 (8/28/2025)

To start off, let's define a line segment between two points $x_1$ and $x_2$. The line segment can be expressed as
\[
    x = \theta x_1 + (1 - \theta) x_2, \quad (\theta \in \mathbb{R})
\]

\begin{figure}[hbpt!]
    \vspace{0.25cm}
    \centering
    \begin{tikzpicture}[scale=2]

        % Draw the main line
        \draw[thick, black] (-2,1) -- (0,0);

        % Extend the line lightly
        \draw[gray!50] (-2.5,1.25) -- (0.5,-0.25);

        \fill (-2,1) circle(0.02) node[above right, yshift=-0.1cm] {$x_1$};
        \fill (0,0) circle(0.02) node[above right, yshift=-0.1cm] {$x_2$};
        \node[below left] at (-2,1) {$\theta = 1$};
        \node[below left] at (0,0) {$\theta = 0$};

        % add points and labels for theta = 1.2, 0.6, -0.2, considering the line is decreasing and x_1 (theta = 1) is at (-2, 1)
        \fill (-2.4,1.2) circle(0.02) node[below left, yshift=-0.1cm] {$\theta = 1.2$};
        \fill (-1.2,0.6) circle(0.02) node[below left, yshift=-0.1cm] {$\theta = 0.6$};
        \fill (0.4,-0.2) circle(0.02) node[below left, yshift=-0.1cm] {$\theta = -0.2$};

    \end{tikzpicture}
\end{figure}
This figure will be useful to visualize the concepts of affine and convex sets. hey!!!!!

\section{Affine set}
We start this chapter by defining an \textit{affine set}. Say we have some $x_1$ and $x_2$. An affine set contains all points on the line through $x_1$ and $x_2$.
\[
    x = \theta x_1 + (1 - \theta) x_2, \quad (\theta \in \mathbb{R})
\]


\textit{Affine set}: contains the line through any two distinct points in the set. And example of this sould
be the solution set of linear equations $\{x \mid Ax + b\}$

\section{Convex set}
\textit{Convex set}: line segment between $x_1$ and $x_2$: all points
\[
    x = \theta x_1 + (1 - \theta) x_2, \quad (0 \leq \theta \leq 1)
\]
The convex set contains the line segment between \textit{any} two points in set.
\[
    x_1, x_2 \in C, \quad 0 \leq \theta \leq 1 \implies \theta x_1 + (1 - \theta) x_2 \in C
\]
Exercise 1

Let $C \subset \mathbb{R}^n$ be a convex set, $x_1$, $x_2$, $\dots$, $x_k \in C$. Let $\theta_1$, $\theta_2$, $\dots$, $\theta_k \in \mathbb{R}$ satisfy $\theta_i \geq 0$ and $\sum_{i=1}^k \theta_i = 1$. Show that

Show that $\theta_1 x_1 + \theta_2 x_2 + \dots + \theta_k x_k \in C$.
\begin{align*}
     & \theta_1 x_1 + \theta_2 x_2 + \dots + \theta_k x_k = 1                                                                                     \\
     & = (\theta_1 + \theta_2) \left( \frac{\theta_1}{\theta_1 + \theta _2} x_1 + \frac{\theta_2}{\theta_1 + \theta_2} x_2 \right) + \theta_3 x_3
\end{align*}
Convex combination and convex hull

Convex combination of $x_1$, $\dots$, $x_k$: any point $x$ of form
\[
    x = \theta_1 x_1 + \theta_2 x_2 + \dots + \theta_k x_k, \quad \theta_i \geq 0, \sum_{i=1}^k \theta_i = 1
\]
Convex hull ($\conv{S}$): set of all convex combinations of points in $S$.


% -- END LECTURE 2 (8/28/2025)

